\subsection{One correction through the elliptic involution and its branch quotient}

\begin{result}{E.188}{Proved}
{The three-torsion correction extends to the actual eight-point quotient}
Use the normalized full-GT action of E.182 and the actual
torus with labelled marks \(T[3]\), as in E.184.
Put \(\iota(z)=-z\), and add the other three points of
\(T[2]\).  Thus the twelve-point configuration is
\(T[2]\cup T[3]\).
The pointwise translation correction of E.184 can be extended
to one correction at twelve points which fixes the named
\(\iota\), induces that same corrected action after forgetting
the three new marks, and fixes every old label.
It may permute the three new labels by one permutation
\(\sigma_2\in\mathfrak S_3\).

Let \(C=C_{\Gamma_{1,[12]}}(\iota)\), and let \(D\) be the
ordinary sphere mapping class group preserving the partition
of its eight marks into the four branch marks \(B\) and the
four unramified pair marks \(J\).  Then the actual quotient
map and its full completion give
\[
 1\longrightarrow\langle\iota\rangle\longrightarrow C
      \longrightarrow D\longrightarrow1,\qquad
 C_{\widehat\Gamma_{1,[12]}}(\iota)=\widehat C,\qquad
 \widehat C/\langle\iota\rangle=\widehat D.
\tag{E.188.1}
\]
In particular, if \(C^0\) is the inverse image of the
pure eight-point group, the corrected action induces an
automorphism of the actual full group
\[
              \widehat C^0/\langle\iota\rangle
                              =\widehat\Gamma_{0,8}.
\tag{E.188.2}
\]
This is an equality for the full induced profinite topologies.
The induced action on the colored label quotient
\(\mathfrak S_B\times\mathfrak S_J\) is conjugation by the
same \(\sigma_2\), extended to fix the original origin
and every member of \(J\).
\end{result}

\begin{proof}
Write \(\alpha_9\) for the action already corrected in E.184,
so that \(\alpha_9(h)=h\) for every named \(h\in H=T[3]\).
The involution \(\iota_9\) lies in the zero-label subgroup
\(\widehat{\mathcal N}_0\) of that result: its conjugation on
\(H\) is \(h\mapsto-h\).  The same is therefore true of
\(\alpha_9(\iota_9)\).
Forgetting the other eight marks identifies
\(\widehat{\mathcal N}_0\) with \(\widehat\Gamma_{1,1}\).
The image of \(\iota_9\) is the central elliptic involution
\(-I\).  The normalized action fixes this element:
in the flat-pants definition it fixes every pants
decomposition, so its weight is one.  Conjugating that
action does not change this central element.
The injectivity of the stated forgetting map consequently gives
\[
                         \alpha_9(\iota_9)=\iota_9.
\tag{E.188.3}
\]

Choose the twelve-point pants marking compatibly with the
nine-point marking.  E.182 supplies the ordinary forgetting
square, including its completed version.  The pure forgetting
map is surjective: its ordinary kernel is the fundamental
group of the ordered configuration of the three additional
points in the punctured torus, and this configuration is
connected.  Lift the one pure conjugator of E.184 through
this surjection.  After this lift, the twelve-point
automorphism induces exactly \(\alpha_9\); denote it for
the moment by \(\widetilde\alpha_{12}\).

There is an exact relative group in which its remaining
involution discrepancy lives.  Set
\[
 Y=\operatorname{Conf}_3(T\setminus T[3]),\qquad
 K=\pi_1(Y),\qquad E_{\rm rel}=K\rtimes\langle\iota_{12}\rangle.
\tag{E.188.4}
\]
The torus and all nine old marks are fixed in this fibre.
Its free extension ranks are \(10,11,12\), and its real
dimension is six.  The involution acts on \(Y\) by
negating each of its three coordinates.  Its fixed set
consists of exactly six points, the six ordered assignments
of \(T[2]\setminus\{0\}\):
\[
                         |Y^\iota|=3!=6.
\tag{E.188.5}
\]
The relative kernel has its full topology in the ambient
group.  Indeed apply the free-normal-kernel argument of
E.182 successively to forgetting the three new marks:
at each stage the completed kernel is the actual
completed punctured-torus free group.  Comparing this
sequence with the corresponding free extension sequence
for \(K\), induction gives the injection of
\(\widehat K\) and identifies its image with the kernel
of the composite forgetting map.  Adjoining the named
finite involution retains this identification.
Equation E.188.3 and the unchanged label action imply that
\(\widetilde\alpha_{12}(\iota_{12})\) is a complement
inside the full completion of E.188.4.

Apply the one-cyclic-pass case of the actual finite-cover
proof E.176 to \(Y\).  The proof of E.176.5 applies to its
three free kernels; the bounded Tate and reduced-cone
argument uses one involution, rather than two.  Finite
fusion and compactness therefore conjugate this complement
to one of the six ordinary complements associated to
E.188.5, by one element of \(\widehat K\).
An isotopy of the torus relative to the nine old points
can move the three additional points along any chosen path
in \(Y\) to the assigned ordering.  Its endpoint gives an
ordinary mapping class permuting only those three new
labels; after forgetting them its mapping class is the
identity.  Its conjugation transports the corresponding
ordinary complement to \(\langle\iota_{12}\rangle\).
Combining these two corrections gives an element whose
nine-point forgetting is the identity and whose only
label permutation is one \(\sigma_2\in\mathfrak S_3\).
This proves the first assertion.  In particular no second,
independently chosen nine-point correction has entered.

We next prove the ordinary centralizer quotient, including
the group that is actually completed.  Removing all
twelve marks makes the involution action free.  The
quotient is the sphere with its eight indicated marks
removed.  Its surface groups give the actual index-two
extension
\[
 1\longrightarrow \Pi_{13}\longrightarrow \Phi_7
       \xrightarrow{\epsilon}C_2\longrightarrow1.
\tag{E.188.6}
\]
Here \(\Pi_{13}\) and \(\Phi_7\) are free of ranks thirteen
and seven, respectively.  The character \(\epsilon\)
is one on the four branch peripherals and zero on the
four pair peripherals.  The relation among all eight
peripherals is consistent with this character because
there are four branch marks.

The conjugation embedding of \(\Phi_7\) in
\(\operatorname{Aut}(\Pi_{13})\) has image precisely the
inverse image of its outer deck group.  Its kernel is
trivial: it meets the center-free finite-index subgroup
\(\Pi_{13}\) trivially and is a finite subgroup of a free
group.  A mapping class centralizing \(\iota\) thus extends
to an outer automorphism of \(\Phi_7\).
The ordinary free-group root property, applied to the
squares of branch peripherals, shows that this extension
preserves every positive peripheral type, with the branch
and pair sets retained separately.
The punctured Dehn--Nielsen--Baer theorem then gives an
ordinary positive mapping class of the quotient sphere;
the genus-one and genus-zero surfaces both have negative
Euler characteristic, so Theorem 8.8 of
\citet[version 5.0, pp.245--246]{LitFarbMargalit2011Version5}
applies in both places.

The kernel consists exactly of the two deck classes.
Indeed if the extended outer automorphism is inner,
its restriction differs from a deck class by an inner
automorphism of \(\Pi_{13}\); ordinary faithfulness on
the punctured torus makes the residual mapping class
the identity.  Conversely, every colored sphere mapping
class preserves \(\epsilon\), lifts through the actual
cover, and its lift commutes with the deck involution.
This proves the first exact sequence in E.188.1.
It also proves surjectivity onto the colored permutation
group \(\mathfrak S_B\times\mathfrak S_J\).

We verify the full induced topology of \(C\), instead of
deducing it from that ordinary exact sequence.
The genus-zero congruence theorem says that
\[
 \widehat D\longrightarrow
          \operatorname{Out}(\widehat\Phi_7)
                          \quad\text{is injective}.
\tag{E.188.7}
\]
It applies to \(D\) because \(D\) is an actual finite-index
subgroup of the eight-punctured sphere mapping class group.
One primary proof of the needed genus-zero theorem is
\citet[\S4, pp.9--10, and Theorem 18, pp.16--18,
in the author version]{LitKent2016}.
More explicitly, the exact profinite Birman sequence and
the injective completed point-pushing map in that proof
propagate injectivity from the pure three-punctured sphere,
whose group is trivial.  Passing from the pure group to
its finite label extension retains injectivity, since the
peripheral permutation is detected in the surface action.
Thus the theorem here concerns full profinite completion,
not merely the image called the procongruence group.

Complete the ordinary exact sequence.  The involution
still embeds, since its permutation exchanges the members
of four marked pairs.  Right exactness gives
\(\widehat C/\langle\iota\rangle=\widehat D\).
Suppose that \(z\in\widehat C\) has trivial outer action
on \(\widehat\Pi_{13}\).  The image of \(z\) in
\(\operatorname{Out}(\widehat\Phi_7)\) restricts, modulo
inner automorphisms of \(\widehat\Pi_{13}\), to a deck
class.  This formulation is necessary: changing a
representative by an inner automorphism of
\(\widehat\Phi_7\) can change the outer restriction by
that deck class.
Choose such an inner change to remove the deck class,
and then an inner change from \(\widehat\Pi_{13}\)
to make the restriction literally the identity.
An automorphism of a finitely generated free profinite
group which is the identity on an open subgroup is the
identity on the whole group, by
\citet[Lemma 2.8, pp.333--334]{LitBuxErshovRapinchuk2011}.
The image of \(z\) in \(\operatorname{Out}(\widehat\Phi_7)\)
is therefore trivial.  Equation E.188.7 makes its image
in \(\widehat D\) trivial, so \(z\in\langle\iota\rangle\).
The nonidentity deck class acts nontrivially on the
puncture conjugacy classes of \(\widehat\Pi_{13}\).
Thus \(z=1\).
We have proved that \(\widehat C\) injects in the outer
surface action, hence also in
\(\widehat\Gamma_{1,[12]}\).  Its topology is full, and
the same is true of its finite-index subgroup \(C^0\).

It remains to identify the entire ambient centralizer
with this embedded completion.  Let \(P\) be the pure
level-six subgroup of \(\Gamma_{1,[12]}\).
Use the actual space \(X\) of eleven ordered nonzero
distinct points on the full-level-six universal torus.
It is a smooth aspherical manifold, with a poly-free
fundamental group \(P\), and real dimension twenty-four.
The finite group \(\Gamma_{1,[12]}/P\) acts on it by
changing the level and label markings.  Its actual
orbifold path extension is \(\Gamma_{1,[12]}\).
The one-pass version of E.176.2, now in this full
finite path extension, shows for every finite output
that every refined centralizer element projects to an
ordinary element of \(C\).  Consequently
\[
 C_{\widehat\Gamma_{1,[12]}}(\iota)=\overline C
                                              =\widehat C.
\tag{E.188.8}
\]
The last equality is the full-topology result just proved.
The finite-cover argument is essential here; no equality
between a discrete fixed locus and a profinite centralizer
has been assumed.

Finally, the corrected automorphism fixes \(\iota\), so
preserves the centralizer E.188.8.  A member of \(C^0\)
has a label permutation that fixes each of the four fixed
marks and each of the four paired orbits, while allowing
exchanges inside a pair.  This condition is preserved by
conjugation of the label quotient by \(\sigma_2\).
Hence the corrected action preserves \(\widehat C^0\),
descends through its named deck subgroup, and induces
E.188.2.  Its colored label quotient is exactly the
permutation claimed in the statement.  All corrections
were made before any of the four pair-forgetting maps
was applied.
\end{proof}
